% Part: incompleteness
% Chapter: representability-in-q
% Section: comp-representable

\documentclass[../../../include/open-logic-section]{subfiles}

\begin{document}

\olfileid{inc}{req}{crq}
\olsection{గణనీయ ప్రమేయాలు~$\Th{Q}$లో ప్రాతినిధ్యం చేయదగినవి}

\begin{thm}
ప్రతి గణనీయ ప్రమేయమూ~$\Th{Q}$లో ప్రాతినిధ్యం చేయదగినది.
\end{thm}

\begin{proof}
నిర్దిష్టత కోసం, చర్చ్--ట్యూరింగ్ సిద్ధాంతప్రతిపాదనను వాడుతూ, ఒక ప్రమేయం
సామాన్య పునరావృత్త ప్రమేయమైతే, అప్పుడు మాత్రమే అది గణనీయమని తీసుకుందాం.
సామాన్య పునరావృత్త ప్రమేయాలు అంటే శూన్య ప్రమేయం~$\Zero$, ఉత్తరగామి
ప్రమేయం~$\Succ$, ప్రక్షేప ప్రమేయం~$\Proj{n}{i}$ల నుంచి సంయుక్తం, ఆదిమ
పునరావృత్తి, సక్రమ కనిష్ఠీకరణలను వాడి నిర్వచించగలిగిన ప్రమేయాలు.
\olref[pri]{lem:prim-rec} వల్ల, $f$, $g$ల నుంచి నిర్వచించగల ఏ ప్రమేయం~$h$నైనా
$f$, $g$, $\Zero$, $\Succ$, $\Proj{n}{i}$, $\Add$, $\Mult$, $\Char{=}$ల
నుంచి సంయుక్తం, సక్రమ కనిష్ఠీకరణలతో నిర్వచించవచ్చు. అందువల్ల ఒక ప్రమేయం
సామాన్య పునరావృత్త ప్రమేయమైతే, అప్పుడు మాత్రమే దాన్ని $\Zero$, $\Succ$,
$\Proj{n}{i}$, $\Add$, $\Mult$, $\Char{=}$ల నుంచి సంయుక్తం, సక్రమ
కనిష్ఠీకరణలతో నిర్వచించవచ్చు.

సంబంధిత ప్రాథమిక ప్రమేయాలు~$\Th{Q}$లో ప్రాతినిధ్యం చేయదగినవని కూడా చూపాం
(\cref{inc:req:bre:prop:rep-zero,inc:req:bre:prop:rep-succ,inc:req:bre:prop:rep-proj,inc:req:bre:prop:rep-id,inc:req:bre:prop:rep-add,inc:req:bre:prop:rep-mult});
ప్రాతినిధ్యం చేయదగిన ప్రమేయాల నుంచి సంయుక్తం లేదా సక్రమ కనిష్ఠీకరణతో
నిర్వచించిన ఏ ప్రమేయమూ
(\olref[cmp]{prop:rep-composition},
\olref[min]{prop:rep-minimization}) ప్రాతినిధ్యం చేయదగినదని కూడా చూపాం.
కాబట్టి ప్రతి సామాన్య పునరావృత్త ప్రమేయమూ~$\Th{Q}$లో ప్రాతినిధ్యం
చేయదగినది.
\end{proof}

\begin{explain}
గణనీయ ప్రమేయాల సమితిని~$\Th{Q}$లో ప్రాతినిధ్యం చేయదగిన ప్రమేయాల సమితిగా
లక్షణీకరించవచ్చని చూపాం. నిజానికి ఈ నిరూపణ మరింత సాధారణమైనది.
ప్రాతినిధ్యయోగ్యత నిర్వచనం నుంచి, $\Th{Q}$ను విస్తరించే (లేదా $\Th{Q}$ను
అర్థనిర్దేశం చేయగల) ఏ సిద్ధాంతమూ గణనీయ ప్రమేయాలకు ప్రాతినిధ్యం చేయగలదని
సులభంగా తెలుస్తుంది. తిరుగుగా, \tetoken{వ్యుత్పత్తి}{!!{derivation}} భావన గణనీయమైన ఏ
\tetoken{వ్యుత్పత్తి}{!!{derivation}} వ్యవస్థలోనైనా, ప్రాతినిధ్యం చేయదగిన ప్రతి ప్రమేయమూ గణనీయమైనదే.
అందువల్ల, ఉదాహరణకు, గణనీయ ప్రమేయాల సమితిని పియానో అంకగణితంలో, లేదా
జెర్మెలో--ఫ్రెంకెల్ సమితి సిద్ధాంతంలో కూడా ప్రాతినిధ్యం చేయదగిన ప్రమేయాల
సమితిగా లక్షణీకరించవచ్చు. గ్యోడెల్ గమనించినట్లుగా ఇది కొంత ఆశ్చర్యకరం.
నిరూప్యత విషయానికి వస్తే, ఏ సిద్ధాంతాన్ని పరిగణిస్తామనే దానిపై ప్రశ్నలు
చాలా సున్నితంగా ఆధారపడతాయని చూస్తాం; స్థూలంగా చెప్పాలంటే, స్వీకృతాలు ఎంత
బలంగా ఉంటే అంత ఎక్కువను నిరూపించవచ్చు. కానీ విస్తృతమైన స్వీకృత సిద్ధాంతాల
పరిధిలో ప్రాతినిధ్యం చేయదగిన ప్రమేయాలు సరిగ్గా గణనీయ ప్రమేయాలే; సిద్ధాంతాలు
స్వీకృతీకరించదగినవిగా ఉన్నంతకాలం, బలమైన సిద్ధాంతాలు మరిన్ని ప్రమేయాలకు
ప్రాతినిధ్యం చేయవు.
\end{explain}

\end{document}
